\begin{textsample}{NEG Dim 4 – 2021\_c4 – Score 1.00 – 2021\_c4/enade\_2021\_tec\_red\_comp\_4311331117.txt}  \label{ex:f4_neg_223}
CAPÍTULO 4 PERCEPÇÃO DA PROVA Nas análises feitas neste capítulo, constam as percepções dos concluintes da Área de Tecnologia em Redes de Computadores sobre a prova aplicada no Enade/2021.
Essas percepções foram mensuradas por meio de nove questões que avaliaram desde o grau de dificuldade da prova até o tempo gasto para concluí -la.
As percepções sobre a prova foram relacionadas ao desempenho dos estudantes, à Grande Região de funcionamento do curso, à Categoria Administrativa e ao tipo de Organização Acadêmica da IES.
O questionário de percepção da prova encontra -se ao final do Anexo VIII, onde está a reprodução do exame.
O desempenho dos estudantes foi classificado em quatro quartos.
Para tanto, esse desempenho foi ordenado de forma ascendente.
O percentil 25, P25, também conhecido como primeiro quartil, é a nota de desempenho que deixa um quarto ( 25 \% ) dos valores observados abaixo e três quartos acima.
Na Figura 1, apresenta -se uma ilustração desse conceito.
O quarto inferior de desempenho é composto pelas notas abaixo do primeiro quartil.
Já o percentil 75, P75, também conhecido como terceiro quartil, é o valor para o qual há três quartos ( 75 \% ) dos dados abaixo e um quarto acima dele.
O quarto superior de desempenho é composto pelas notas iguais ou acima do terceiro quartil.
O percentil 50, P50, também conhecido como mediana, é o valor que divide as notas em dois conjuntos de igual tamanho.
O segundo quarto inclui os valores entre o primeiro quartil ( P25 ) e a mediana.
O terceiro quarto contém os valores entre a mediana ( P50 ) e o terceiro quartil ( P75 ).
Vale ressaltar que percentis, quartis e medianas são pontos que, não obrigatoriamente, pertencem ao conjunto original de dados, ao passo que os quartos são subconjuntos dos dados originais.
A seguir, serão apresentados gráficos com resultados selecionados, relativos às nove questões avaliadas por grupos de estudantes.
Nas barras dos gráficos, apresenta -se o percentual de alunos que assinalou uma das opções ou a soma das porcentagens daqueles que assinalaram duas ( ou três ) delas.
Por exemplo, para as questões 1 e 2, nos gráficos, é apresentada a porcentagem total de participantes que assinalou as opções ( D ) Difícil e ( E ) Muito difícil.
Em cada barra, foram assinalados, também, os extremos do intervalo de confiança de 95 \% como linhas verticais unidas por uma linha horizontal na forma da letra H maiúscula.
O estimador de um parâmetro com um certo nível de confiança ( por exemplo, 95,0 \% ) deve conter o parâmetro no intervalo de confiança em 95 \% das vezes.
Na comparação entre os estimadores dos parâmetros de duas classes de uma dada categoria ( por exemplo : Norte e Nordeste nas Grandes Regiões, de primeiro e último quarto dentro de desempenho, de IES Públicas e Privada, ou de Universidades e Faculdades ), associados aos seus respectivos intervalos de confiança, diz -se que não há diferença estatisticamente significativa entre os parâmetros das duas categorias quando há interseção entre os intervalos de confiança, e que há diferença, se os intervalos de confiança são disjuntos ( para maiores informações vide Glossário ).
Nas Tabelas do Anexo II, são apresentados os valores absolutos e a distribuição percentual¹⁸ das alternativas válidas das nove questões, segundo o mesmo recorte de desempenho dos alunos e da Grande Região de funcionamento do curso, de Categoria Administrativa e do tipo de Organização Acadêmica da IES. 4.1 GRAU DE DIFICULDADE DA PROVA 4.1.1 COMPONENTE DE FORMAÇÃO GERAL Ao avaliarem “ Qual o grau de dificuldade desta prova na parte de Formação Geral? ” ( Questão 1 ), 42,0 \% do grupo de inscritos e presentes optaram pelas alternativas Difícil ou Muito difícil.
Entretanto, para uma maior parte dos estudantes ( 49,0 \% ), o Componente de Formação Geral da prova foi avaliado com grau de dificuldade Médio ( Gráfico 4.1, Gráfico 4.2 e, no Anexo II, a Tabela II.1 ).
O percentual de estudantes que consideraram a prova Difícil ou Muito difícil foi maior na região Norte, onde a proporção foi de 51,0 \%, enquanto a de menor proporção foi a Nordeste, com 34,4 \%.
No Gráfico 4.1, é possível observar que as diferenças entre a região Nordeste e as regiões Sudeste e Sul são estatisticamente significativas.
Nas Grandes Regiões, a proporção de presentes à prova que avaliaram o Componente de Formação Geral como de grau de dificuldade Médio esteve entre 45,1 \%, na região Norte, e 56,2 \%, na região Nordeste. como de grau de dificuldade Médio esteve entre 45,1 \%, na região Norte, e 56,2 \%, na região Nordeste.
O percentual de alunos que consideraram o componente de Formação Geral da prova Difícil ou Muito difícil foi decrescente em função dos quartos de desempenho : 44,7 \%, no primeiro quarto, e 39,8 \%, no quarto, o grupo de melhor desempenho na prova.
Nos quartos de desempenho intermediários, a proporção de alunos que consideraram a prova Difícil ou Muito difícil foi de 42,3 \%, no segundo quarto, e de 41,3 \%, no terceiro quarto.
As diferenças entre os quartos de desempenho não são estatisticamente significativas.
Para todos os quartos de desempenho, a alternativa modal a essa pergunta foi a Médio, com 46,1 \% e 51,1 \% dos respondentes nos quartos extremos, primeiro e quarto, respectivamente.
O percentual de estudantes que consideraram o componente de Formação Geral da prova Difícil ou Muito difícil foi maior entre os estudantes de Universidades, alcançando 45,0 \%, enquanto entre os estudantes das Centros Universitários, a Organização Acadêmica com a menor incidência, alcançou 36,4 \%.
No Gráfico 4.3, é possível observar que a diferença entre Universidades e Centros Universitários é estatisticamente significativa.
Nas Organizações Acadêmicas, a proporção de presentes à prova que avaliaram o Componente de Formação Geral como de grau de dificuldade Médio esteve entre 46,5 \%, nas Universidades, e 55,1 \% nos Centros Universitários ( ver também a Tabela II.2, no Anexo II ).
O percentual de estudantes que consideraram o componente de Formação Geral da prova Difícil ou Muito difícil foi maior entre os estudantes de cursos de IES Privadas ( 43,1 \% ), sem diferença estatisticamente significativa da proporção registrada em IES Públicas ( 38,1 \% ).
Já a proporção de presentes à prova que avaliaram o Componente de Formação Geral como de grau de dificuldade Médio foi de 52,1 \%, entre os estudantes de IES Públicas, e de 48,2 \%, entre os de IES Privadas ( ver Gráfico 4.4 ). 4.1.2 COMPONENTE DE CONHECIMENTO ESPECÍFICO Ao responderem à Questão 2 – “ Qual o grau de dificuldade desta prova na parte de Componente Específico? ” – 60,2 \% do grupo de estudantes classificaram -na como Difícil ou Muito difícil.
O Componente modal de Conhecimento Específico da prova foi o grau Difícil, com quase a metade dos estudantes ( 49,5 \% ) assim avaliando seu nível de dificuldade.
Além disso, mais de um terço dos alunos ( 36,1 \% ) considerou Médio o grau de dificuldade da prova ( Gráfico 4.5, Gráfico 4.6, e, no Anexo II, a Tabela II.3 ).
A análise das respostas dos estudantes quanto ao grau de dificuldade do Componente de Conhecimento Específico da prova, agregado por Grande Região, não mostra diferenças estatisticamente significativas entre proporções de alunos que a avaliaram como Difícil ou Muito difícil.
A maior proporção foi encontrada na região Norte ( 68,6 \% ), e a menor, na região Nordeste ( 54,4 \% ).
Nas demais regiões, as proporções foram : 61,1 \%, na região Sudeste ; 63,4 \%, na região Sul e 60,6 \%, na região Centro-Oeste.
O percentual de alunos que classificaram o grau de dificuldade como Médio, no Componente de Conhecimento Específico, variou de 31,4 \% a 42,1 \%, nas regiões Norte e Nordeste, respectivamente.
Considerando -se a avaliação da dificuldade das questões do Componente de Conhecimento Específico da prova de acordo com o desempenho dos estudantes, observa -se uma tendência crescente dos que consideraram a parte específica Difícil ou Muito difícil do primeiro ao terceiro quarto de desempenho, variando de 58,1 \% ( primeiro quarto ) a 64,4 \% ( terceiro quarto ).
No quarto superior de desempenho ( 59,3 \% ), a proporção decai, chegando ao patamar do segundo quarto de desempenho ( 59,1 \% ).
Observa -se que não há diferença estatisticamente significativa dos resultados entre os quartos de desempenho.
Já os que responderam que o grau de dificuldade das questões do Componente de Conhecimento Específico da prova foi Médio atingiram percentagens que variaram de 33,0 \%, no terceiro quarto, a 38,3 \%, no último.
Considerando -se a avaliação da dificuldade das questões do Componente de Conhecimento Específico da prova, o percentual de estudantes que as consideraram Difícil ou Muito difícil foi maior entre os estudantes de Universidades, atingindo 62,6 \% ; já entre os de Faculdades, a Organização Acadêmica com a menor incidência, essa proporção foi de 55,5 \%.
No Gráfico 4.7, é possível observar que não há diferença estatisticamente significativa entre os quatro tipos de Organização Acadêmica.
Nas Organizações Acadêmicas, a proporção de presentes à prova que consideraram o Componente de Conhecimento Específico de grau de dificuldade Médio esteve entre 33,8 \%, nas Universidades, e 41,5 \%, nas Faculdades ( ver também Tabela II.4, no Anexo II ).
O percentual de estudantes que consideraram o Componente Específico da prova Difícil ou Muito difícil foi maior entre os estudantes de cursos de IES Privadas ( 60,9 \% ), sem diferença estatisticamente significativa da proporção registrada em IES Públicas ( 57,7 \% ).
Já a proporção de presentes à prova que avaliaram esse componente como de grau de dificuldade Médio foi de 40,2 \%, entre os estudantes de IES Públicas, e de 35,0 \%, entre os de IES Privadas ( ver Gráfico 4.8 e Tabela II.4, no Anexo II ). 4.2 EXTENSÃO DA PROVA EM RELAÇÃO AO TEMPO TOTAL Indagados quanto à extensão da prova, em relação ao tempo total oferecido para a sua resolução ( Questão 3 ), os estudantes apontaram, com maior incidência, a alternativa que considerava a extensão Adequada, em todas as agregações consideradas ( Gráfico 4.9, Gráfico 4.10, e, no Anexo II, a Tabela II.5 ).
O percentual de alunos que responderam ser a extensão da prova Adequada foi de 62,8 \%.
Já 30,6 \% dos inscritos presentes consideraram que a prova foi Longa ou Muito longa, e 6,6 \% a avaliaram como Curta ou Muito curta.
Entre as Grandes Regiões, os que avaliaram a prova como Longa ou Muito longa, em relação ao tempo total destinado à sua resolução, atingiram uma porcentagem que variou de 21,6 \%, na região Norte, até 39,7 \%, na região Sul.
Há diferença estatisticamente significativa entre a região Sul e as regiões Norte, Sudeste e Centro-Oeste.
Considerando -se o desempenho dos alunos, nota -se ainda que 64,6 \% consideraram a extensão da prova Adequada, no quarto de desempenho inferior, e 61,2 \%, no de melhor desempenho ( quarto superior ).
Nos quartos intermediários, essa proporção foi de 62,6 \%, no segundo quarto, e de 62,7 \%, no terceiro.
No Gráfico 4.10, observa -se que as proporções de estudantes que consideraram a prova Longa ou Muito longa apresentam tendência crescente em função dos quartos de desempenho, sendo a menor proporção dessa resposta no primeiro quarto de desempenho ( 28,1 \% ), e a maior, no terceiro quarto ( 32,1 \% ).
Pode -se constatar que não há diferença estatisticamente significativa entre os quartos de desempenho.
Considerando -se a extensão da prova, o percentual de estudantes que avaliaram a prova como Longa ou Muito longa, em relação ao tempo total destinado à sua resolução, foi maior entre os estudantes de Faculdades, alcançando 32,1 \%, enquanto entre os estudantes de CEFET/IF, a Organização Acadêmica com a menor incidência, esse índice atingiu 27,4 \% ; já entre os de Centros Universitários, o índice atingiu 29,7 \%, e entre os de Universidades, 31,0 \%.
No Gráfico 4.11, é possível observar que não há diferença estatisticamente significativa entre os quatro tipos de Organização Acadêmica.
A proporção de presentes à prova que consideraram a sua extensão Adequada, em relação ao tempo total destinado à sua resolução, esteve entre 62,0 \%, nas Faculdades, e 63,9 \%, nos Centros Universitários ( ver também Tabela II.6, no Anexo II ).
O percentual de estudantes que consideraram a prova Longa ou Muito longa, em relação ao tempo total destinado à sua resolução, foi maior entre os estudantes de cursos de IES Privadas ( 31,7 \% ), não havendo diferença estatisticamente significativa da proporção em relação aos de IES Públicas ( 26,5 \% ).
Já a proporção de presentes à prova que consideraram a extensão Adequada foi de 64,4 \%, entre os estudantes de IES Públicas, e de 62,3 \%, entre os de IES Privadas ( ver Gráfico 4.12 e Tabela II.6, no Anexo II ). 4.3 COMPREENSÃO DOS ENUNCIADOS DAS QUESTÕES 4.3.1 COMPONENTE DE FORMAÇÃO GERAL Em relação aos enunciados das questões do Componente de Formação Geral ( Questão 4 ), as opiniões foram positivas, já que 70,3 \% dos alunos avaliados consideraram todos ou a maioria dos enunciados das questões “ claros e objetivos ” ( Gráfico 4.13, Gráfico 4.14 e, no Anexo II, a Tabela II.7 ).
Na análise regional, a percentagem de estudantes que avaliaram que todos ou a maioria dos enunciados das questões do Componente de Formação Geral estavam “ claros e objetivos ” variou de 62,7 \%, na região Norte, a 74,5 \%, na região Sul.
A análise das percepções dos estudantes sobre a clareza e a objetividade dos enunciados permite afirmar que todos, ou a maioria dos enunciados de questões relativas ao Componente de Formação Geral, foram considerados “ claros e objetivos ” para a maior parte dos respondentes ( maior ou igual a 62,7 \%, em todas as regiões, e maior ou igual a 66,1 \%, em todos os quartos de desempenho ).
Segundo o desempenho, observa -se que a proporção dos que emitiram essa opinião cresce, conforme o desempenho aumenta, com diferença estatisticamente significativa entre o quarto inferior de desempenho e o quarto mais elevado.
No quarto superior, a clareza e a objetividade de todos ou da maioria dos enunciados das questões foram percebidas por 74,4 \%.
A proporção dos presentes que consideraram que todos ou a maioria dos enunciados das questões do Componente de Formação Geral estavam “ claros e objetivos ” foi maior entre os estudantes de Centros Universitários e CEFET/IF, chegando a 70,9 \% em ambos, enquanto entre os de Faculdades, a Organização Acadêmica com a menor incidência, esse índice atingiu 69,3 \% ; já entre os de Universidades, o índice foi de 70,3 \%.
No Gráfico 4.15, é possível observar que as diferenças entre os quatro tipos de Organização Acadêmica não são estatisticamente significativas ( ver também Tabela II.8, no Anexo II ).
O percentual de estudantes que consideraram que todos ou a maioria dos enunciados das questões do Componente de Formação Geral estavam “ claros e objetivos ” foi maior entre os estudantes de cursos de IES Privadas ( 70,6 \% ), mas sem diferença estatisticamente significativa da proporção registrada em IES Públicas ( 69,2 \% ).
Um maior detalhamento pode ser encontrado no Gráfico 4.16 e na Tabela II.8, no Anexo II. 4.3.2 COMPONENTE DE CONHECIMENTO ESPECÍFICO Em relação aos enunciados das questões do Componente de Conhecimento Específico da prova, para 72,4 \% dos estudantes avaliados da Área de Tecnologia em Redes de Computadores, a clareza e a objetividade ( Questão 5 ) estavam presentes em todas ou na maioria das questões ( Gráfico 4.17 a Gráfico 4.20, e no Anexo II, as Tabelas II.9 e II.10 ).
A maioria dos estudantes de todas as Grandes Regiões brasileiras considerou “ claros e objetivos ” todos ou a maioria dos enunciados das questões do Componente de Conhecimento Específico da prova, percentual sempre maior ou igual a 62,7 \%.
A diferença entre a região Sudeste e a região Sul é estatisticamente significativa.
A proporção de estudantes que consideraram os enunciados das questões “ claros e objetivos ” apresenta uma tendência crescente em relação ao aumento de desempenho : mais elevada no quarto superior ( 80,5 \% ), se comparada ao quarto inferior de desempenho ( 66,2 \% ).
As diferenças entre o último quarto de desempenho e os dois primeiros quartos são estatisticamente significativas, assim como entre o primeiro e o terceiro quartos.
A proporção dos presentes que consideraram que todos ou a maioria dos enunciados das questões do Componente de Conhecimento Específico estavam “ claros e objetivos ” foi maior entre os estudantes de CEFET/IF, atingindo 76,8 \%.
Entre os estudantes de Faculdades, a Organização Acadêmica com a menor incidência, essa proporção atingiu 67,8 \%, sendo intermediárias as proporções registradas em Centro Universitários ( 75,2 \% ) e Universidades ( 72,3 \% ).
No Gráfico 4.19, é possível observar que as diferenças entre os tipos de Organização Acadêmica não são estatisticamente significativas ( ver também Tabela II.10, no Anexo II ).
O percentual de estudantes que consideraram que todos ou a maioria dos enunciados das questões do Componente de Conhecimento Específico estavam “ claros e objetivos ” foi maior entre os estudantes de cursos de IES Públicas ( 73,3 \% ), sem diferença estatisticamente significativa da proporção registrada em IES Privadas ( 72,2 \% ).
Ver Gráfico 4.20 e Tabela II.10, no Anexo II, para maior detalhamento. 4.4 SUFICIÊNCIA DAS INFORMAÇÕES/INSTRUÇÕES FORNECIDAS Ao se avaliarem as informações/instruções fornecidas para a resolução das questões ( Questão 6 ), constatou -se que 79,2 \% dos respondentes da Área de Tecnologia em Redes de Computadores de todo o Brasil afirmaram que eram até excessivas ou suficientes em todas ou na maioria das questões ( Gráfico 4.21, Gráfico 4.22 e, no Anexo II, a Tabela II.11 ).
Quanto à distribuição de respondentes pelas Grandes Regiões, observa -se que a proporção de estudantes que consideraram as informações/instruções fornecidas até excessivas ou suficientes em todas ou na maioria das questões foi sempre superior ou igual a 72,5 \%, chegando a 83,9 \%, na região Sul.
Não existem diferenças estatisticamente significativas entre as Grandes Regiões.
Levando -se em conta o desempenho dos participantes, observa -se uma tendência, grosso modo, crescente das proporções de participantes que avaliaram as informações/instruções como até excessivas ou suficientes em todas ou na maioria das questões, havendo diferenças estatisticamente significativas entre o primeiro quarto e os dois quartos superiores, bem como entre o segundo e o terceiro quartos de desempenho, conforme mostrado no Gráfico 4.22.
O percentual foi mais elevado no terceiro quarto ( 83,9 \% ), percentual superior à média nacional ( 79,2 \% ).
Já no quarto inferior, a suficiência das informações/instruções declarada como até excessiva, em todas ou na maioria das questões foi percebida por 73,5 \% dos respondentes.
Considerando -se as informações/instruções fornecidas para a resolução das questões ( Questão 6 ), a proporção que considerou que eram até excessivas ou suficientes em todas ou na maioria das questões foi maior entre os estudantes de CEFET/IF, totalizando 82,3 \%.
As Universidades foram a Organização Acadêmica em que houve a menor proporção : 78,3 \%.
No Gráfico 4.23, não é possível observar diferenças estatisticamente significativas entre os quatro tipos de Organização Acadêmica ( ver também Tabela II.12, no Anexo II ).
O percentual de estudantes que consideraram que as informações/instruções fornecidas para a resolução das questões eram até excessivas ou suficientes em todas ou na maioria das questões foi maior entre os estudantes de cursos de IES Públicas ( 80,8 \% ), não havendo diferença estatisticamente significativa da proporção registrada em IES Privadas ( 78,7 \% ) ( ver Gráfico 4.24 e Tabela II.12, no Anexo II ). 4.5 DIFICULDADE ENCONTRADA AO RESPONDER À PROVA Indagados sobre as dificuldades com as quais se depararam ao responder à prova ( Questão 7 ), 21,9 \% dos estudantes apontaram o Desconhecimento do conteúdo.
Para 47,6 \%, a Forma diferente de abordagem do conteúdo foi indicada como dificuldade.
Já a Falta de motivação para fazer a prova foi a dificuldade apontada por 16,3 \% dos respondentes.
Considerando -se todo o Brasil, 10,6 \% dos respondentes afirmaram que não tiveram qualquer tipo de dificuldade para responder à prova ( Tabelas II.13 e II.14, no Anexo II ).
Nos Gráficos 4.25 a 4.28, são apresentados os percentuais de estudantes que apontaram o Desconhecimento do conteúdo como “ dificuldade ao responder à prova ”.
Na análise por Grandes Regiões, o percentual de inscritos e presentes que apontaram o Desconhecimento do conteúdo como “ dificuldade ao responder à prova ” não superou 25,5 \%.
Os percentuais variaram de 20,3 \%, na região Centro-Oeste, a 25,5 \%, na região Norte.
A Forma diferente de abordagem do conteúdo foi a escolha modal dos estudantes, com percentuais que variaram de 46,9 \% ( região Nordeste ) a 52,3 \% ( região Centro-Oeste ).
O percentual de alunos que citaram a Falta de motivação como dificuldade variou de 12,5 \% ( região Centro-Oeste ) a 17,7 \% ( região Nordeste ).
Os que declararam não ter qualquer tipo de dificuldade para responder à prova variaram de 5,9 \%, na região Norte, a 13,1 \%, na região Sul.
Em relação aos quartos de desempenho, o Desconhecimento do conteúdo foi a opção escolhida por 18,3 \% dos estudantes do quarto inferior e por 28,3 \% do quarto superior, caracterizando uma tendência crescente.
São observadas diferenças estatisticamente significativas entre o quarto superior e os demais quartos de desempenho.
A alternativa modal entre os alunos, quando agregados pelos quartos de desempenho, para a dificuldade encontrada, foi causada pela Forma diferente de abordagem do conteúdo : 44,7 \%, no último quarto, e 50,0 \%, no terceiro quarto, escolheram essa alternativa.
Na análise por Organização Acadêmica, o percentual de inscritos e presentes que apontou o Desconhecimento do conteúdo como “ dificuldade ao responder à prova ” não superou 25,3 \%.
Os percentuais variaram de 19,2 \%, nos Centros Universitários, a 25,3 \%, nas Faculdades.
A Forma diferente de abordagem do conteúdo foi a escolha modal dos estudantes, com percentuais que variaram de 41,4 \% ( Faculdades ) a 51,4 \% ( CEFET/IF ).
O percentual de alunos que citaram a Falta de motivação como dificuldade variou de 15,5 \% ( Universidades ) a 17,4 \% ( Faculdades ).
Os que declararam não ter qualquer tipo de dificuldade para responder à prova variaram de 5,9 \%, nas CEFET/IF, a 11,6 \%, nos Universidades.
No Gráfico 4.27, é possível observar que não há diferenças estatisticamente significativas entre os quatro tipos de Organização Acadêmica ( ver também Tabela II.14, no Anexo II ).
Em relação às Categorias Administrativas, o Desconhecimento do conteúdo foi a opção escolhida por 24,1 \% dos estudantes, nas IES Públicas, e por 21,3 \%, nas IES Privadas.
Não são observadas diferenças estatisticamente significativas entre essas duas categorias.
A alternativa modal entre os alunos, quando agregados pelas Categorias Administrativas, para a dificuldade encontrada, foi causada pela Forma diferente de abordagem do conteúdo : 44,9 \%, entre os estudantes de IES Públicas, e 48,3 \%, entre os de IES Privadas ( ver Gráfico 4.28 e Tabela II.14, no Anexo II ). 4.6 CONTEÚDOS DAS QUESTÕES OBJETIVAS DA PROVA Ao se analisarem os conteúdos das questões objetivas da prova ( Questão 8 ), verifica -se que um percentual pequeno dos estudantes avaliados, apenas 10,9 \%, afirmou que não estudou ainda a maioria desses conteúdos ( Gráficos 4.29 a Gráfico 4.32, e as Tabelas II.15 e II.16, no Anexo II ).
A maior parte dos estudantes ( 43,6 \% ) afirmou ter estudado e aprendido muitos ou todos os conteúdos avaliados.
Na análise por Grande Região, a proporção de respondentes que escolheram a opção não estudou ainda a maioria desses conteúdos foi pequena.
Observa -se que, nas regiões Norte ( 15,7 \% ), Sudeste ( 11,0 \% ) e Centro-Oeste ( 20,3 \% ), as proporções foram maiores do que a média nacional ( 10,9 \% ).
Nas regiões Nordeste e Sul, a situação foi inversa, com proporções mais baixas, correspondendo a, respectivamente, 9,3 \% e 7,2 \%.
São observadas diferenças estatisticamente significativas entre a região Centro-Oeste e as regiões Nordeste, Sudeste e Sul.
Em todas as regiões, a maioria dos presentes afirmou ter estudado e aprendido muitos ou todos os conteúdos, com proporções variando entre 33,3 \%, na região Norte, a 50,3 \%, na região Nordeste.
Considerando -se separadamente as opiniões de estudantes dos quatro quartos de desempenho, observa -se que, no quarto inferior, 16,8 \% ofereceram como resposta que não estudou ainda a maioria desses conteúdos, sendo 6,1 \% os do quarto superior com a mesma resposta.
A diferença entre os alunos que optaram por esse motivo de dificuldade no quarto inferior em relação aos demais quartos de desempenho é estatisticamente significativa.
Tendo -se em conta o quarto superior, 56,0 \% dos alunos afirmaram ter estudado e aprendido muitos ou todos os conteúdos.
No outro extremo, no primeiro quarto, 33,7 \% optaram pelas mesmas categorias.
As proporções são crescentes com o desempenho.
Na análise por Organização Acadêmica, a proporção de respondentes que escolheram a opção não estudou ainda a maioria desses conteúdos foi pequena.
Observa -se que a proporção nas Faculdades, 9,3 \%, foi menor do que a da média nacional ( 10,9 \% ).
Nas Universidades, Centros Universitários e CEFET/IF, a situação foi inversa : a proporção foi mais alta ( 11,3 \%, 11,4 \% e 11,0 \%, respectivamente ).
No Gráfico 4.31, é possível observar que as diferenças entre os quatro tipos de Organização Acadêmica não são estatisticamente significativas.
Em todos os tipos de Organização Acadêmica, a maioria dos presentes afirmou ter estudado e aprendido muitos ou todos os conteúdos, com proporções variando entre 41,8 \%, nas Universidades, a 49,3 \%, nos CEFET/IF ( ver também Tabela II.16, no Anexo II ).
Em relação às Categorias Administrativas, observa -se que, nas IES Públicas, 9,8 \% ofereceram como resposta que não estudou ainda a maioria desses conteúdos, sendo 11,2 \% nas IES Privadas com a mesma resposta.
A diferença entre os alunos que optaram por esse motivo de dificuldade nas duas categorias não é estatisticamente significativa.
Considerando -se as IES Públicas, 47,5 \% dos alunos afirmaram ter estudado e aprendido muitos ou todos os conteúdos, sendo de 42,5 \% a proporção dos que escolheram essa resposta nas IES Privadas ( ver Gráfico 4.32 e Tabela II.16, no Anexo II ). 4.7 TEMPO GASTO PARA CONCLUIR A PROVA Ao responderem sobre o tempo de conclusão da prova ( Questão 9 ), mais de dois terços dos estudantes ( 68,4 \% ) afirmaram ter gasto Entre duas e quatro horas ( Gráfico 4.33 a Gráfico 4.36 e, no Anexo II, as Tabelas II.17 e II.18 ).
Considerando -se as cinco Grandes Regiões brasileiras, os percentuais dos que utilizaram Entre duas e quatro horas para finalizar a prova na região Sudeste ( 65,3 \% ) foi inferior ao percentual nacional.
Nas demais Grandes Regiões, os percentuais de alunos que dispensaram Entre duas e quatro horas para concluir a prova foram mais elevados que a média : 80,0 \%, na região Norte ; 76,0 \%, na região Nordeste ; 68,7 \%, na região Sul ; e 69,8 \%, na região Centro-Oeste, como mostrado no Gráfico 4.33.
Há diferenças estatisticamente significativas entre a região Sudeste e as regiões Norte e Nordeste.
Uma vez analisadas às alternativas escolhidas pelos estudantes que se situam nos diferentes quartos de desempenho, observa -se uma tendência crescente : uma proporção maior de participantes no quarto superior declarou ter gasto Entre duas e quatro horas para concluir a prova quando comparada com a proporção de participantes nos quartos inferiores.
Há diferenças estatisticamente significativas do quarto inferior para o terceiro e o último quartos de desempenho, bem como entre o segundo e o quarto superior de desempenho.
Na análise por Organização Acadêmica, os percentuais dos que utilizaram Entre duas e quatro horas para finalizar a prova nos Centros Universitários ( 71,0 \% ) e nos CEFET/IF ( 78,8 \% ) foram superiores ao percentual nacional ( 68,4 \% ).
Nas Faculdades e Universidades, tais percentuais foram de 65,7 \% e 66,5 \%, respectivamente.
Como mostrado no Gráfico 4.35, há diferença estatisticamente significativa dos CEFET/IF em relação às Universidades e às Faculdades.
Em relação às Categorias Administrativas, observa -se que, nas IES Públicas, 71,2 \% declararam ter gasto Entre duas e quatro horas para concluir a prova, sendo 67,7 \% nas IES Privadas com a mesma resposta.
A diferença entre os alunos que optaram por esse motivo de dificuldade nas duas categorias não é estatisticamente significativa ( ver Gráfico 4.36 e Tabela II.18, no Anexo II ).

% matched lemmas: 
\end{textsample}
